%% Crystal structure prediction with a geometry-aware coordinate loss.
%%
%% WORKING DOCUMENT. At this stage it records measured results only; the
%% introduction, related work, method and discussion are deliberately absent.
%% Numbers here are copied from topn_summary_<TAG>.csv and nothing else.

\documentclass[manuscript,screen,review]{acmart}

\AtBeginDocument{%
  \providecommand\BibTeX{{%
    Bib\TeX}}}

%% Suppressed while this is an internal record rather than a submission.
\settopmatter{printacmref=false}
\setcopyright{none}
\renewcommand\footnotetextcopyrightpermission[1]{}
\pagestyle{plain}

\begin{document}

\title{A Periodicity-Aware Coordinate Loss for Autoregressive Crystal
Structure Prediction}

%% TODO: author name and affiliation.
\author{Anonymous}
\affiliation{%
  \institution{University of South Carolina}
  \city{Columbia}
  \state{South Carolina}
  \country{USA}}
\email{dg47@email.sc.edu}

\renewcommand{\shortauthors}{Working draft}

\begin{abstract}
  Autoregressive crystal generators discretize fractional coordinates into
  bins and train with cross-entropy over those bins. That objective is
  geometrically blind: it penalizes a one-bin error as heavily as a half-cell
  error, and it treats the adjacent coordinates $0.99$ and $0.01$ as maximally
  distant. We replace it with a circular soft-target loss, in which probability
  mass is distributed over neighbouring bins under a wrapped Gaussian whose
  distance metric respects the periodicity of the unit cell. Evaluating on a
  180-structure crystal structure prediction benchmark, against training data
  held disjoint from it, we find that the modified objective does not improve on
  plain cross-entropy: at kernel widths of 2 and 64 bins it is equal or worse on
  structure match and substantially worse on space-group match, and the penalty
  grows with the width of the kernel rather than appearing at one particular
  scale. We attribute this to soft targets blurring each coordinate
  independently, where symmetry depends on exact relations between coordinates,
  and conclude that the per-coordinate factorization of the output, rather than
  the shape of the target placed on it, is what limits symmetry recovery.
\end{abstract}

\keywords{crystal structure prediction, autoregressive models, discretized
coordinates, periodic loss functions}

\maketitle

\section{Results}
\label{sec:results}

\subsection{Experimental setup}
\label{sec:setup}

Every model here is a LLaMA-style autoregressive transformer
(512-dimensional, $4.3\times10^{7}$ parameters) over crystal token sequences of
the form \texttt{[SOS] [ATOMS] (element$\vert$ox, $x$, $y$, $z$)$^{N}$}
\texttt{[LATTICE] ($a,b,c,\alpha,\beta,\gamma$) [EOS]}, with fractional
coordinates and lattice parameters quantized to 1024 bins. Oxidation states are
assigned by composition-based guessing. Training runs for 500 epochs, and the
checkpoint with the lowest validation loss is the one evaluated.
Section~\ref{sec:changes} states what distinguishes the individual runs.

All runs train on the same split of MP-20: 27{,}023 training, 9{,}001
validation and 9{,}009 test structures, held disjoint from the benchmark by
excluding every structure whose reduced composition matches one of the 180
benchmark formulas. Matching on the \emph{reduced} composition makes
Zn$_4$S$_4$O$_{16}$, ZnSO$_4$ and a $2\times$ supercell collide, as it should,
since the task supplies a formula and asks for a structure. Every rate reported
below is therefore an out-of-sample measurement.

Evaluation uses the 180 benchmark structures. For each formula, 50 structures are
generated with the composition constrained to the requested formula, relaxed
with MatRIS (\texttt{matris\_10m\_oam}, FIRE, \texttt{FrechetCellFilter}, at
most 200 steps, $f_{\max} = 0.05$\,eV/\AA), and ranked by formation energy per
atom. Match rate at top-$N$ scores only the $N$ lowest-energy candidates.
Structure match uses \texttt{StructureMatcher} with
$\mathrm{ltol} = 0.2$, $\mathrm{stol} = 0.3$ and $\mathrm{angle\_tol} = 5$;
space groups are assigned with $\mathrm{symprec} = 0.1$. Consensus requires a
single candidate to match both structure and space group. All rates are over
the full set of 180 formulas: a formula that produced no valid structure is
counted as a failure rather than excluded.

\subsection{The original setup, before any changes}
\label{sec:original}

The starting point was the Materium code as released. It treats a crystal as a
sentence. Each atom contributes four tokens --- which element it is, and its
three fractional coordinates --- and the unit cell contributes six more at the
end. Fractional coordinates are not stored as numbers but rounded into one of
1024 equally spaced bins, so a coordinate of $0.2504$ becomes "bin 256". The
model is an ordinary autoregressive transformer: it reads the tokens produced so
far and predicts the next one, and it is trained with plain cross-entropy, the
same objective used for predicting the next word in a sentence.

Figure~\ref{fig:pipeline} shows how a prediction is produced and scored. Nothing
in that path was changed by us except where Section~\ref{sec:changes} says so.

\begin{figure}[h]
\begin{footnotesize}
\begin{verbatim}
  formula  e.g. Zn4S4O16
     |
     v
  GENERATE   transformer emits tokens left to right
     |       50 candidates per formula
     |       composition is held fixed during sampling
     v
  RELAX      MatRIS, 200 steps, machine-learned forces
     |       gives each candidate a formation energy
     v
  RANK       sort by energy, keep the best N
     |
     v
  MATCH      compare against the known structure
             pymatgen StructureMatcher
             report: structure / space group / both
\end{verbatim}
\end{footnotesize}
\caption{One formula in, a ranked list of structures out. Match rate at top-$N$
scores only the $N$ lowest-energy candidates.}
\label{fig:pipeline}
\end{figure}

Two properties of that original objective matter. First, it is blind to
geometry: predicting bin 257 when the answer is bin 256 is punished exactly as
hard as predicting bin 800. Second, it does not know the cell repeats, so bins
1023 and 0 are treated as far apart even though they describe neighbouring
positions. Both observations motivated the loss work reported here.

\subsection{What changed in each run}
\label{sec:changes}

The three completed runs share their training data, their bin grid, their
oxidation handling and their epoch budget, and none of them uses augmentation or
an auxiliary output. They differ in the training objective alone, which is what
makes Table~\ref{tab:changes} short and the comparison in
Section~\ref{sec:completed} a controlled one.

\begin{table}[h]
  \caption{What distinguishes each completed run. The kernel width $\sigma$ is
  given in bins and, in parentheses, in fractional coordinates. The evaluated
  checkpoint is the one with the lowest validation loss.}
  \label{tab:changes}
  \begin{footnotesize}
  \begin{tabular}{llll}
    \toprule
    Run & Loss & Kernel width $\sigma$ & Eval.\ epoch \\
    \midrule
    C1 & flat          & ---                 & 159 \\
    C2 & soft, narrow  & 2 bins ($0.002$)    & 311 \\
    C3 & soft, wide    & 64 bins ($0.063$)   & 251 \\
    \bottomrule
  \end{tabular}
  \end{footnotesize}
\end{table}

In plain terms:

\paragraph{Loss: flat versus soft.}
\emph{Flat} is the original cross-entropy. \emph{Soft} spreads the training
target over neighbouring bins instead of demanding one exact bin, and treats bin
1023 and bin 0 as adjacent. A near miss is therefore rewarded, and the
periodicity of the cell is respected.

\paragraph{Bin grid.}
The original code converted a coordinate to a bin by multiplying by 1023 but
converted back by dividing by 1024. Encoding and decoding sat on grids offset by
about one and a half bin widths, so a structure could not be written and read
back unchanged. The corrected grid uses the same spacing in both directions,
and all three runs below use it. This is a correctness fix, not a modelling
choice, and it is the one respect in which these runs differ from the released
code regardless of their objective.

\paragraph{Augmentation.}
Where a crystal's origin sits is arbitrary: sliding every atom by the same
amount describes the same material. Augmentation exploits that by shifting all
coordinates by a random amount each time a structure is shown to the model, so
it sees many equivalent versions. We later found this was implemented
incorrectly --- it reproduced the ordering the tokenizer would emit for a
shifted structure only $31\%$ of the time --- and disabled it for every later
run.

\subsection{Completed results}
\label{sec:completed}

Table~\ref{tab:completed} gives the runs finished so far. All rates are over
the full 180 benchmark structures, with a formula that produced nothing counted
as a failure. One structure is $0.56$ percentage points, so a gap of one point
is under two structures.

\begin{table}[h]
  \caption{Match rate (\%) on the 180-structure benchmark. Each block is one
  run; rows within it are the number of lowest-energy candidates scored.
  Str = structure, SG = space group, Con = both.}
  \label{tab:completed}
  \begin{footnotesize}
  \begin{tabular}{lrrrrrrrrrrrr}
    \toprule
    & \multicolumn{3}{c}{top-1} & \multicolumn{3}{c}{top-5}
    & \multicolumn{3}{c}{top-10} & \multicolumn{3}{c}{top-25} \\
    \cmidrule(lr){2-4}\cmidrule(lr){5-7}\cmidrule(lr){8-10}\cmidrule(lr){11-13}
    Run & Str & SG & Con & Str & SG & Con & Str & SG & Con & Str & SG & Con \\
    \midrule
    C1 & 48.3 & 40.6 & 38.9 & 52.8 & 48.9 & 46.7 & 53.9 & 50.0 & 47.2 & 54.4 & 51.7 & 48.9 \\
    C2 & 47.8 & 38.3 & 37.2 & 51.1 & 43.3 & 41.7 & 53.3 & 45.6 & 44.4 & 53.9 & 48.9 & 47.8 \\
    C3 & 45.0 & 33.9 & 32.2 & 48.3 & 36.7 & 35.0 & 51.1 & 38.9 & 36.7 & 52.2 & 40.0 & 37.8 \\
    \bottomrule
  \end{tabular}
  \end{footnotesize}
\end{table}

The ordering is the same at every budget, and it is monotone in the width of
the target kernel.

\paragraph{Flat cross-entropy is the strongest objective tested.}
C1 leads on all three metrics at every budget but one. At top-25 it reaches
$54.4$ on structure, $51.7$ on space group and $48.9$ on both together, against
$53.9 / 48.9 / 47.8$ for the narrow kernel and $52.2 / 40.0 / 37.8$ for the wide
one. The original objective, in other words, is not the thing holding these
models back.

\paragraph{Widening the kernel costs symmetry, not geometry.}
Across the whole range from flat to $\sigma = 64$, the structure rate at top-25
moves by $2.2$ points while the space-group rate moves by $11.7$ and the
consensus rate by $11.1$. The damage is concentrated almost entirely in
symmetry, and it grows with $\sigma$ rather than appearing at one particular
width. That rules out the most hopeful reading of the narrow result --- that
$\sigma = 2$ was simply too small to reach the scale at which the model errs,
and that a wider kernel would do better. A wider kernel does worse, steadily.

\paragraph{Why the mechanism predicts exactly this.}
A soft target spreads probability over neighbouring bins, so a near miss on any
one coordinate is cheap. Symmetry, however, is a property of \emph{exact
relations} between coordinates --- an atom at $0.25$ paired with one at $0.75$,
or an atom sitting on a special position. Nudging each coordinate independently
costs almost nothing under a wrapped-Gaussian target but destroys those
relations, and destroying them is precisely what the space-group test detects.
Mean positional accuracy and symmetry are not the same quantity, and this
objective buys the first at the expense of the second.

\paragraph{What cannot be read off it.}
The three runs stopped at different epochs, because the soft-loss variants
ceased improving earlier than flat cross-entropy; C2's and C3's evaluated
checkpoints are their own best, not an arbitrary truncation, but the comparison
is not at equal optimisation effort. Differences under about a point correspond
to fewer than two structures and should not be read as effects. Nothing here
separates the choice of objective from the choice of \emph{representation}: all
three runs predict independent per-coordinate bins, and that is the shared
assumption the space-group gap most likely indicts.

\section{Status and open items}
\label{sec:status}

Recorded here so that measured results and pending work stay separable.

\begin{itemize}
  \item The runs in Table~\ref{tab:completed} stopped at different epochs,
        because the soft-loss variants ceased improving earlier than flat
        cross-entropy. Differences of under about a point correspond to fewer
        than two structures out of 180 and should not be read as effects.
  \item The kernel-width question is settled for this family of objectives.
        At $\sigma = 2$ bins the kernel is far narrower than the measured
        validation coordinate error of $\approx 0.125$, which suggested a wider
        kernel should help; at $\sigma = 64$ it is substantially worse on every
        metric. Further variation of $\sigma$ is not worth the compute.
  \item An auxiliary space-group output, trained alongside the tokens, is
        under evaluation. Its purpose is to force symmetry into the internal
        representation rather than into the coordinate targets.
  \item Translation augmentation is disabled everywhere. It reproduced the
        tokenizer's own site ordering for a shifted structure only $31\%$ of
        the time, so it was training the model on sequences the tokenizer would
        never produce. The cause is not yet identified; the cell reduction
        applied before tokenization is the main suspect, since that is also not
        invariant to relabelling the cell axes.
  \item Objectives that change what the model can represent --- the auxiliary
        space-group head, a loss on interatomic distances under the
        minimum-image convention, or a Wyckoff-based representation that
        enforces symmetry by construction --- address the measured deficit more
        directly than further variation of the coordinate kernel.
\end{itemize}

\end{document}
